\documentclass[10pt,a4paper]{amsart}
\usepackage{amsmath,amssymb,amsthm,graphicx}
\usepackage{verbatim}
\newenvironment{reference}{\comment}{\endcomment}
%\newenvironment{reference}{}{}
\newenvironment{slogan}{\comment}{\endcomment}
\newenvironment{history}{\comment}{\endcomment}

% For commutative diagrams we use Xy-pic
\usepackage[all]{xy}

% We use 2cell for 2-commutative diagrams.
\xyoption{2cell}
% Theorem environments.
%
\theoremstyle{plain}
\newtheorem{theorem}[subsection]{Theorem}
\newtheorem{proposition}[subsection]{Proposition}
\newtheorem{lemma}[subsection]{Lemma}

\theoremstyle{definition}
\newtheorem{definition}[subsection]{Definition}
\newtheorem{example}[subsection]{Example}
\newtheorem{exercise}[subsection]{Exercise}
\newtheorem{situation}[subsection]{Situation}

\theoremstyle{remark}
\newtheorem{remark}[subsection]{Remark}
\newtheorem{remarks}[subsection]{Remarks}

\numberwithin{equation}{subsection}

% Macros
%
\def\lim{\mathop{\mathrm{lim}}\nolimits}
\def\colim{\mathop{\mathrm{colim}}\nolimits}
\def\Spec{\mathop{\mathrm{Spec}}}
\def\Hom{\mathop{\mathrm{Hom}}\nolimits}
\def\Ext{\mathop{\mathrm{Ext}}\nolimits}
\def\SheafHom{\mathop{\mathcal{H}\!\mathit{om}}\nolimits}
\def\SheafExt{\mathop{\mathcal{E}\!\mathit{xt}}\nolimits}
\def\Sch{\mathit{Sch}}
\def\Mor{\mathop{\mathrm{Mor}}\nolimits}
\def\Ob{\mathop{\mathrm{Ob}}\nolimits}
\def\Sh{\mathop{\mathit{Sh}}\nolimits}
\def\NL{\mathop{N\!L}\nolimits}
\def\CH{\mathop{\mathrm{CH}}\nolimits}
\def\proetale{{pro\text{-}\acute{e}tale}}
\def\etale{{\acute{e}tale}}
\def\QCoh{\mathit{QCoh}}
\def\Ker{\mathop{\mathrm{Ker}}}
\def\Im{\mathop{\mathrm{Im}}}
\def\Coker{\mathop{\mathrm{Coker}}}
\def\Coim{\mathop{\mathrm{Coim}}}

% Boxtimes
%
\DeclareMathSymbol{\boxtimes}{\mathbin}{AMSa}{"02}

%
% Macros for moduli stacks/spaces
%
\def\QCohstack{\mathcal{QC}\!\mathit{oh}}
\def\Cohstack{\mathcal{C}\!\mathit{oh}}
\def\Spacesstack{\mathcal{S}\!\mathit{paces}}
\def\Quotfunctor{\mathrm{Quot}}
\def\Hilbfunctor{\mathrm{Hilb}}
\def\Curvesstack{\mathcal{C}\!\mathit{urves}}
\def\Polarizedstack{\mathcal{P}\!\mathit{olarized}}
\def\Complexesstack{\mathcal{C}\!\mathit{omplexes}}
% \Pic is the operator that assigns to X its picard group, usage \Pic(X)
% \Picardstack_{X/B} denotes the Picard stack of X over B
% \Picardfunctor_{X/B} denotes the Picard functor of X over B
\def\Pic{\mathop{\mathrm{Pic}}\nolimits}
\def\Picardstack{\mathcal{P}\!\mathit{ic}}
\def\Picardfunctor{\mathrm{Pic}}
\def\Deformationcategory{\mathcal{D}\!\mathit{ef}}
\usepackage{cleveref}
\crefname{originallemma}{original lemma}{original lemmas}
\Crefname{originallemma}{Original lemma}{Original lemmas}
\makeatletter
\newlabel{lemma-sieve-sheafification}{{00ZO}{source}{}{sourcepointer-1}{}}
\newlabel{lemma-sieve-sheafification@cref}{{[originallemma][0][]00ZO}{[1][0][]source}{}{}{}}
\makeatother
\title{Comparing two topologies by inclusion of their sheaves}
\author{The Stacks Project Authors}
\date{stacks-project:sites.tex:local-snapshot-20261012; GNU FDL 1.2 or later}
\begin{document}
\maketitle
\noindent\textit{Editorial note.} All delivered mathematical text is the author's own, copied from the pinned Stacks Project chapter file sites-cohomology.tex; no statement or proof step was rewritten, reordered or removed.\par\medskip
\noindent\textit{Editorial note.} The proof is the author's complete argument; the phrase 'By formal nonsense' introduces the construction of the comparison map between the two sheafifications, which the following lines describe explicitly, and it is preserved verbatim inside the counted span.\par\medskip
\noindent\textit{Editorial note.} Every labelled reference occurring inside the counted spans is carried in assembly.external\_references as a non-load-bearing navigation pointer (same-chapter ones at their definition site, cross-chapter ones at the invocation inside the counted span), so no load-bearing mathematics is replaced by a stub. No bibliography citation occurs inside the counted spans.\par\medskip
\noindent\textit{Editorial note.} The checked chapter list for same-name collisions is the set of stacks chapters staged in this repository: chow, cohomology, divisors, etale-cohomology, flat, more-algebra, more-morphisms, morphisms, spaces-more-morphisms, varieties together with this card's chapter sites-cohomology.\par\medskip
\noindent\textit{Editorial note.} The preamble is the chapter's own retained minimal preamble (sources/stacks-sites-cohomology/retained-preamble.tex), a verbatim excerpt of the native preamble: the Xy-pic setup and the native theorem declarations with the mathematical macros.\par\medskip
\noindent\textit{Allowed standard prerequisites (editorial).}\begin{itemize}
\item Covering sieves characterized by sheafification. Same chapter, Lemma lemma-sieve-sheafification. Supplies the criterion applied to the sieve generated by an arbitrary covering.
\item Sheafification as a functor determined by the category of sheaves. Same chapter, Proposition proposition-sheafification-adjoint and Theorem theorem-topology-and-topos. Allows the comparison of the two sheafifications induced by the inclusion of the categories of sheaves.
\end{itemize}\medskip
\ifdefined\hypertarget\hypertarget{sourcepointer-1}{}\fi
\noindent\textit{Source reference.} Original lemma 00ZO: sources/stacks-sites/pinned-original-source.tex, source lines 11965--11965. Navigation pointer invoked at proof line 11965; the counted span contains the pointer only and the cited result is not reproduced inside it, so the pointer is registered as a non-load-bearing navigation stub in the established form for Stacks material. No bibliography citation is involved.\par\medskip
\begin{theorem}
\label{theorem-topology-and-topos}
Let $\mathcal{C}$ be a category.
Let $J$, $J'$ be topologies on $\mathcal{C}$.
The following are equivalent
\begin{enumerate}
\item $J = J'$,
\item sheaves for the topology $J$ are the same as
sheaves for the topology $J'$.
\end{enumerate}
\end{theorem}

\begin{lemma}
\label{lemma-finer-topology}
Assumption and notation as in Theorem \ref{theorem-topology-and-topos}.
Then $J \subset J'$ if and only if every sheaf for the
topology $J'$ is a sheaf for the topology $J$.
\end{lemma}

\begin{proof}
One direction is clear. For the other direction suppose that
$\Sh(\mathcal{C}, J') \subset \Sh(\mathcal{C}, J)$.
By formal nonsense this implies
that if $\mathcal{F}$ is a presheaf of sets,
and $\mathcal{F} \to \mathcal{F}^\#$,
resp.\ $\mathcal{F} \to \mathcal{F}^{\#, \prime}$
is the sheafification wrt $J$, resp.\ $J'$ then there
is a canonical map $\mathcal{F}^\# \to \mathcal{F}^{\#, \prime}$
such that
$\mathcal{F} \to \mathcal{F}^\# \to \mathcal{F}^{\#, \prime}$
equals the canonical map $\mathcal{F} \to \mathcal{F}^{\#, \prime}$.
Of course, $\mathcal{F}^\# \to \mathcal{F}^{\#, \prime}$
identifies the second sheaf as the sheafification of the first
with respect to the topology $J'$.
Apply this to the map $S \to h_U$ of
Lemma \ref{lemma-sieve-sheafification}. We get a commutative
diagram
$$
\xymatrix{
S \ar[r] \ar[d] &
S^\# \ar[r] \ar[d] &
S^{\#, \prime} \ar[d] \\
h_U \ar[r] &
h_U^\# \ar[r] &
h_U^{\#, \prime}
}
$$
And clearly, if $S$ is a covering sieve for the topology $J$
then the middle vertical map is an isomorphism (by the lemma)
and we conclude that the right vertical map is an isomorphism as
it is the sheafification of the one in the middle wrt $J'$.
By the lemma again we conclude that $S$ is a covering sieve
for $J'$ as well.
\end{proof}

\end{document}
